\documentclass{standalone}
\usepackage{ctex}
\usepackage{tikz}
\usetikzlibrary{er,calc}

\tikzset{
    entity/.append style = {
        draw,
        thick,
        minimum width=3.0cm,
        minimum height=1.1cm,
        font=\small\bfseries
    },
    relationship/.append style = {
        draw,
        thick,
        aspect=2,
        font=\small\bfseries
    },
    attribute/.append style = {
        draw,
        thick,
        minimum width=2.7cm,
        minimum height=0.9cm,
        font=\scriptsize
    },
    link/.style = {
        thick
    }
}

% #1 实体节点名
% #2 属性节点前缀
% #3 属性半径
% #4 起始角度
% #5 属性列表
\newcommand{\PlaceAttrs}[5]{%
    \def\AttrCount{0}%
    \foreach \txt [count=\k] in {#5} {%
        \xdef\AttrCount{\k}%
    }%
    \pgfmathsetmacro{\AngleStep}{360/\AttrCount}%
    \foreach \txt [count=\k from 1] in {#5} {%
        \pgfmathsetmacro{\AngleNow}{#4+(\k-1)*\AngleStep}%
        \pgfmathsetmacro{\OppAngle}{mod(\AngleNow+180,360)}%
        \node[attribute] (#2\k) at ($(#1)+(\AngleNow:#3)$) {\txt};
        \draw[link] (#1.\AngleNow) -- (#2\k.\OppAngle);
    }%
}

\begin{document}
\begin{tikzpicture}

%==================== 参数区 ====================
% 三组实体统一使用相同的左右水平距离
\def\LeftX{-7.2}
\def\RightX{7.2}

% 三行实体上下距离
\def\TopY{13.2}
\def\MidY{0}
\def\BotY{-13.2}

% 默认属性到实体的距离
\def\AttrRadius{5.7cm}

% 单独控制部分实体的属性距离
\def\PostAttrRadius{6.2cm}
\def\EmailVerificationAttrRadius{4.8cm}
\def\ConfigAttrRadius{4.8cm}

%==================== 第一行：USER 与 EMAIL_VERIFICATION ====================
\node[entity] (user) at (\LeftX,\TopY) {USER};
\node[relationship] (verifyRel) at (0,\TopY) {生成验证码};
\node[entity] (emailv) at (\RightX,\TopY) {EMAIL\_VERIFICATION};

\draw[link] (user) -- node[above,font=\small] {1} (verifyRel);
\draw[link] (verifyRel) -- node[above,font=\small] {N} (emailv);

\PlaceAttrs{user}{u}{\AttrRadius}{115}{
uuid (PK),
username,
email,
password\_hash,
role,
status,
created\_at,
updated\_at,
last\_login\_at,
last\_login\_ip,
email\_verified
}

\PlaceAttrs{emailv}{ev}{\EmailVerificationAttrRadius}{105}{
uuid (PK),
user\_uuid (FK),
type,
code,
expires\_at,
used\_at,
ip,
created\_at
}

%==================== 第二行：POST 与 FRIEND_LINK ====================
\node[entity] (post) at (\LeftX,\MidY) {POST};
\node[entity] (friend) at (\RightX,\MidY) {FRIEND\_LINK};

\PlaceAttrs{post}{p}{\PostAttrRadius}{96}{
uuid (PK),
title,
content,
content\_text,
cover\_image,
status,
tags,
slug,
excerpt,
view\_count,
pinned\_at,
published\_at,
deleted\_at,
created\_at,
updated\_at
}

\PlaceAttrs{friend}{f}{\AttrRadius}{100}{
uuid (PK),
name,
url,
avatar\_url,
description,
tags,
status,
applicant\_email,
reject\_reason,
joined\_at,
created\_at,
updated\_at
}

%==================== 第三行：CONFIG 与 EMAIL_LOG ====================
\node[entity] (config) at (\LeftX,\BotY) {CONFIG};
\node[entity] (elog) at (\RightX,\BotY) {EMAIL\_LOG};

\PlaceAttrs{config}{c}{\ConfigAttrRadius}{110}{
uuid (PK),
config\_key,
config\_value,
description,
is\_public,
created\_at,
updated\_at
}

\PlaceAttrs{elog}{e}{\AttrRadius}{100}{
uuid (PK),
type,
from\_name,
from\_email,
to,
subject,
status,
error\_message,
params,
is\_read,
created\_at
}

\end{tikzpicture}
\end{document}
